\documentclass[11pt]{amsart}
\usepackage{amsmath,amssymb,amsthm}
\usepackage[margin=1.15in]{geometry}
\usepackage[hidelinks]{hyperref}
\newtheorem{theorem}{Theorem}[section]
\newtheorem{corollary}[theorem]{Corollary}
\newtheorem{proposition}[theorem]{Proposition}
\newtheorem{lemma}[theorem]{Lemma}
\theoremstyle{remark}
\newtheorem{remark}[theorem]{Remark}
\newcommand{\E}{\mathbb{E}}
\newcommand{\Prob}{\mathbb{P}}
\newcommand{\Var}{\operatorname{Var}}
\newcommand{\Cov}{\operatorname{Cov}}
\newcommand{\inv}{\operatorname{inv}}
\newcommand{\dW}{d_{\mathrm W}}
\newcommand{\doi}[1]{\href{https://doi.org/#1}{doi:\nolinkurl{#1}}}
\newcommand{\arxiv}[1]{\href{https://arxiv.org/abs/#1}{arXiv:\nolinkurl{#1}}}
\setlength{\emergencystretch}{2.5em}

\begin{document}

\title[Inversions and two-sided descents]{A joint central limit theorem for inversions and two-sided descents}
\author{Yitzchak Shmalo}
\address{Einstein Institute of Mathematics, The Hebrew University of Jerusalem, Jerusalem, Israel}
\email{yspennstate@gmail.com}
\subjclass[2020]{60F05, 05A05, 60C05}
\keywords{random permutation, inversions, descents, inverse descents, central limit theorem, Stein's method}

\begin{abstract}
Let $\pi$ be a uniform random permutation of $n$ elements, let $I(\pi)$ be its number of inversions and let
$T(\pi)=D(\pi)+D(\pi^{-1})$ be the number of descents of $\pi$ plus the number of descents of $\pi^{-1}$. We show that
$I$ and $T$, centred and scaled by their standard deviations, converge jointly to a pair of independent standard normal
variables, with the Wasserstein distance of every unit linear combination to the standard normal law at most a constant
times $n^{-1/2}$. The same holds with $T$ replaced by any statistic that adds a fixed function over all windows of $k$
consecutive positions of $\pi$ and of $\pi^{-1}$, provided its variance grows linearly, and the triple
$(I,D(\pi),D(\pi^{-1}))$, standardized in the same way, converges to a standard normal vector in $\mathbb R^3$. Exactly,
$\Cov(I,T)=(n-1)/2$. The proof replaces $I$ by its Hoeffding projection on the point process that generates $\pi$
and applies Chatterjee's normal approximation by interaction graphs, as Chatterjee and Diaconis did for $T$ alone.
\end{abstract}

\maketitle

\section{Introduction}\label{sec:intro}

For a permutation $\pi$ of $[n]=\{1,\dots,n\}$ write $I(\pi)=\#\{i<j:\pi(i)>\pi(j)\}$ for the number of inversions
and $D(\pi)=\#\{1\le i<n:\pi(i)>\pi(i+1)\}$ for the number of descents. When $\pi$ is uniform, both are
asymptotically normal. For $I$ this is the normality of Kendall's rank statistic~\cite{Kendall}, a U-statistic in the
sense of Hoeffding~\cite{Hoeffding}; for $D$ it follows from the representation of descents as a one-dependent
sequence. Chatterjee and Diaconis~\cite{CD} studied the two-sided statistic
\[
T(\pi)=D(\pi)+D(\pi^{-1}),
\]
the number of descents of $\pi$ plus the number of descents of its inverse, and proved a central limit theorem for
it with a Wasserstein rate $n^{-1/2}$, using the method of interaction graphs introduced by Chatterjee~\cite{Ch08}.
The asymptotic independence of the descent counts of $\pi$ and $\pi^{-1}$, and with it the normality of $T$, had been
obtained earlier by Vatutin~\cite{Vatutin}. The two-sided statistic was later studied on finite Coxeter groups by
Röttger~\cite{Rottger}, Brück and Röttger~\cite{BruckRottger} and Féray~\cite{FerayCoxeter}, and for Mallows
permutations by He~\cite{He}.

The pair $(I,D)$ is also understood. Fang and Röllin~\cite[Corollary~3.8]{FangRollin} proved that the standardized
vector $(I,D)$ is close to a standard normal vector in $\mathbb R^2$, with an error of order $n^{-1/2}$ on convex sets,
and Dörr and Heiny~\cite{DorrHeiny} also proved the joint central limit theorem and studied the joint extremes. Since
$\pi\mapsto\pi^{-1}$ preserves the uniform law and $I(\pi^{-1})=I(\pi)$, the pair $(I(\pi),D(\pi^{-1}))$ has the law of
$(I,D)$, so this result covers the inverse descents as well. It does not settle the joint law of $I$ with $T$, which
involves $D(\pi)$ and $D(\pi^{-1})$ simultaneously: two bivariate normal limits that share a coordinate do not determine
the trivariate limit. In the outlook of his paper on generalized inversions and descents, Dörr~\cite[Section~5]{Dorr}
notes that the joint asymptotic normality of $T$ with the number of inversions has not been investigated. This paper
proves it.

Fix $k\ge2$. For a sequence of $k$ distinct reals let $\operatorname{st}$ denote its pattern, the permutation of $[k]$
with the same relative order. For functions $f,g$ from the symmetric group $S_k$ to $[-1,1]$ set
\begin{equation}\label{eq:W}
F(\pi)=\sum_{i=1}^{n-k+1}f\bigl(\operatorname{st}(\pi(i),\dots,\pi(i+k-1))\bigr),\qquad
G(\pi)=F_g(\pi^{-1}),\qquad W=F(\pi)+G(\pi),
\end{equation}
where $F_g$ is defined as $F$ with $g$ in place of $f$. We call $W$ a two-sided local statistic of window length $k$.
The function applied to a window does not depend on the window's position; this is used in Lemma~\ref{lem:graph}.
The statistic $T$ is the case $k=2$, $f=g=\mathbf 1\{\cdot=21\}$, and $D(\pi)$, $D(\pi^{-1})$, the numbers of peaks
of $\pi$ and of $\pi^{-1}$ are of the same form, and so is a linear combination of statistics with a common window
length once it is scaled so that $|f|,|g|\le1$. For a random variable $X$ with
positive variance write $X'=(X-\E X)/\sqrt{\Var X}$, and let $\dW(X,Y)=\sup\{|\E h(X)-\E h(Y)|:h\text{ is 1-Lipschitz}\}$ be the Wasserstein distance and $Z$ a standard normal
variable.

\begin{theorem}\label{thm:main}
Fix $k\ge2$ and $c>0$. There are constants $C$ and $n_1$, depending only on $k$ and $c$, such that the following holds.
Let $\pi$ be uniform on $S_n$ with $n\ge n_1$, and let $W$ be a two-sided local statistic of window length $k$ with
$\Var W\ge cn$. Then for every $(a,b)\in\mathbb R^2$ with $a^2+b^2=1$,
\[
\dW\bigl(aI'+bW',Z\bigr)\le Cn^{-1/2}.
\]
In particular, if $W=W_n$ is such a statistic for each $n$ with $\liminf_n\Var(W_n)/n>0$, then
$(I',W_n')$ converges in distribution to a standard normal vector in $\mathbb R^2$.
\end{theorem}

The bound is a family of one-dimensional bounds, one for each direction; it is not a bound in a multivariate
distance. The convergence in distribution follows from it by the Cramér--Wold device.

\begin{corollary}\label{cor:T}
For $n\ge2$, $\Cov(I,T)=(n-1)/2$, and $\Var T=(n+7)/6-1/n$. The correlation of $I$ and $T$ is
$(3\sqrt6+o(1))/n$, and $(I',T')$ converges in distribution to a standard normal vector in $\mathbb R^2$, with
$\dW(aI'+bT',Z)\le Cn^{-1/2}$ uniformly over $a^2+b^2=1$ for all $n\ge n_1$, where $C$ and $n_1$ are universal constants.
\end{corollary}

\begin{corollary}\label{cor:triple}
The vector $\bigl(I',D(\pi)',D(\pi^{-1})'\bigr)$ converges in distribution to a standard normal vector in
$\mathbb R^3$.
\end{corollary}

Three observations make the proof work. The first is the coupling of Chatterjee and Diaconis~\cite{CD}: $\pi$ is the
permutation induced by $n$ independent uniform points of the unit square, so that $\pi$ and $\pi^{-1}$ are read from
the same points by the two coordinate orders. On this space $I$ is the number of discordant pairs of points, a
U-statistic of order two, and its Hoeffding projection~\cite{Hoeffding} is a sum of independent terms, one per point
(Lemma~\ref{lem:hoeff}). A sum of functions of single coordinates creates no interactions in the sense
of~\cite{Ch08}, so the interaction graph that Chatterjee and Diaconis built for $T$ serves without change for every
linear combination of the projection with $W$ (Lemma~\ref{lem:graph}). The second is an exact conditioning identity
for consecutive windows (Lemma~\ref{lem:window}), which shows that $I$ and $W$ are almost uncorrelated, so that the
limiting covariance is the identity. The third is that the projection is within $\sqrt{14/(2n+5)}$ of $I'$ in $L^2$, which
transfers the rate from the projection to $I$ itself. Each step is elementary; the joint statement does not follow from
the marginal results or from the bivariate results with $D$ alone.

Section~\ref{sec:moments} proves the covariance identities, Section~\ref{sec:proj} the projection,
Section~\ref{sec:graph} the interaction graph and the replacement bound, and Section~\ref{sec:proofs} the theorem and
its corollaries.

\section{Covariances}\label{sec:moments}

For $\sigma\in S_k$ write $\inv(\sigma)$ for its number of inversions. If $\sigma$ is uniform on $S_k$, then
$\Var\inv(\sigma)=k(k-1)(2k+5)/72$; for $k=n$ this is $\Var I$.

\begin{lemma}\label{lem:window}
Let $2\le k\le n$, let $\pi$ be uniform on $S_n$, let $1\le i\le n-k+1$ and let $\sigma=\operatorname{st}(\pi(i),\dots,\pi(i+k-1))$.
Then
\[
\E\bigl(I(\pi)\mid\sigma\bigr)=\inv(\sigma)+\E I(\pi)-\E\inv(\sigma).
\]
Consequently $\Cov\bigl(I(\pi),h(\sigma)\bigr)=\Cov\bigl(\inv(\sigma),h(\sigma)\bigr)$ for every $h:S_k\to\mathbb R$,
and the same holds with $\pi^{-1}$ in place of $\pi$ in the definition of $\sigma$.
\end{lemma}

\begin{proof}
Let $B=\{i,\dots,i+k-1\}$, let $S=\pi(B)$ be the set of values in the window, and let $\rho$ be the restriction of
$\pi$ to $[n]\setminus B$. The map $\pi\mapsto(\sigma,S,\rho)$ is a bijection onto the set of triples with
$\sigma\in S_k$, $S$ a $k$-subset of $[n]$ and $\rho$ a bijection from $[n]\setminus B$ onto $[n]\setminus S$; hence
$\sigma$ is uniform and independent of $(S,\rho)$. Split the inversions into the pairs inside $B$, the pairs outside
$B$, and the pairs with one position $p\notin B$ and one position $q\in B$. The first count is $\inv(\sigma)$ and the
second is a function of $\rho$. Because $B$ is a set of consecutive positions, every $p\notin B$ lies either to the
left of all of $B$ or to the right of all of $B$. In the first case the number of $q\in B$ forming an inversion with
$p$ is $\#\{s\in S:s<\rho(p)\}$, in the second case it is $\#\{s\in S:s>\rho(p)\}$; neither depends on $\sigma$. So
$I(\pi)=\inv(\sigma)+\Phi(S,\rho)$ for a function $\Phi$, and the independence of $\sigma$ from $(S,\rho)$ gives
$\E(I\mid\sigma)=\inv(\sigma)+\E\Phi$. Taking expectations identifies the constant. The covariance identity follows
from $\Cov(I,h(\sigma))=\Cov(\E(I\mid\sigma),h(\sigma))$. For $\pi^{-1}$ use that $\pi^{-1}$ is uniform and
$I(\pi^{-1})=I(\pi)$.
\end{proof}

The lemma needs the window to consist of consecutive positions: for a non-consecutive set of positions the
cross term depends on how the window's values are arranged relative to the positions in its gaps.

\begin{lemma}\label{lem:cov}
Let $W$ be a two-sided local statistic of window length $k\le n$. Then
\[
\bigl|\Cov(I,W)\bigr|\le 2(n-k+1)\sqrt{k(k-1)(2k+5)/72}\le 2n\sqrt{k^3/24}.
\]
Moreover $\Cov(I,D(\pi))=\Cov(I,D(\pi^{-1}))=(n-1)/4$ and $\Cov(I,T)=(n-1)/2$.
\end{lemma}

\begin{proof}
By Lemma~\ref{lem:window}, each of the $n-k+1$ windows of $\pi$ contributes
$\Cov(\inv(\sigma),f(\sigma))$ to $\Cov(I,F)$, where $\sigma$ is uniform on $S_k$, and by the Cauchy--Schwarz
inequality and $|f|\le1$ this is at most $\sqrt{\Var\inv(\sigma)}$ in absolute value. The windows of $\pi^{-1}$ are
treated in the same way. The second inequality is $k(k-1)(2k+5)\le3k^3$, that is, $k^2-3k+5>0$. For $k=2$ and $f$ the
indicator of $21$, $\inv(\sigma)=f(\sigma)$ is a Bernoulli$(1/2)$ variable, so each window contributes $1/4$.
\end{proof}

The value $\Var T=(n+7)/6-1/n$ for $n\ge2$ is computed in~\cite{CD}. Together with the classical value
$\Var D(\pi)=\Var D(\pi^{-1})=(n+1)/12$ it gives $\Cov(D(\pi),D(\pi^{-1}))=(n-1)/(2n)$. An enumeration of $S_n$ for $n\le8$ confirms these formulas and
Lemma~\ref{lem:cov}.

\section{The point model and the projection of the inversion count}\label{sec:proj}

Let $P_1,\dots,P_n$ be independent uniform points of $[0,1]^2$, $P_a=(U_a,V_a)$. With probability one the $U_a$ are
distinct and the $V_a$ are distinct. List the points by increasing first coordinate and let $\pi(i)$ be the rank, among
$V_1,\dots,V_n$, of the second coordinate of the $i$-th point. Then $\pi$ is uniform on $S_n$, and $\pi^{-1}(j)$ is
the rank of the first coordinate of the point whose second coordinate has rank $j$. A descent of $\pi$ is a pair of
points adjacent in the first coordinate order whose second coordinates decrease; a descent of $\pi^{-1}$ is a pair
adjacent in the second order whose first coordinates decrease. An inversion of $\pi$ is a discordant pair:
\[
I=\sum_{a<b}h(P_a,P_b),\qquad h(p,q)=\mathbf 1\{(u-u')(v-v')<0\}\quad\text{for }p=(u,v),\ q=(u',v').
\]
This is the coupling of~\cite[Section~5]{CD}.

\begin{lemma}\label{lem:hoeff}
Let $n\ge2$, let $\psi(u,v)=-2(u-\tfrac12)(v-\tfrac12)$ and $L=6n^{-1/2}\sum_{a=1}^n\psi(P_a)$. Then $\E L=0$, $\Var L=1$,
$|\psi|\le\tfrac12$, and
\[
\E\bigl(I'-L\bigr)^2\le\frac{14}{2n+5}.
\]
\end{lemma}

\begin{proof}
For $p=(u,v)$ and $P_2$ uniform, $\E h(p,P_2)=u(1-v)+(1-u)v=\tfrac12+\psi(p)$. Hence $\E h=\tfrac12$, $\E\psi=0$ and
$\E\psi^2=4\cdot(1/12)^2=1/36$, so $\Var L=1$. Put $\tilde h(p,q)=h(p,q)-\tfrac12-\psi(p)-\psi(q)$. Then
$\E(\tilde h(P_1,P_2)\mid P_1)=0$, so the variables $\tilde h(P_a,P_b)$, $a<b$, are uncorrelated with each other
(also when two pairs share a point) and with every $\psi(P_c)$, and
$\E\tilde h^2=\Var h-2\Var\psi=\tfrac14-\tfrac2{36}=\tfrac7{36}$. Summing over pairs,
\[
I-\E I=(n-1)\sum_a\psi(P_a)+R,\qquad R=\sum_{a<b}\tilde h(P_a,P_b),\qquad \E R^2=\binom n2\frac7{36}.
\]
This gives $\Var I=n(n-1)^2/36+7n(n-1)/72=n(n-1)(2n+5)/72$. Dividing by $\sqrt{\Var I}$,
$I'=r_nL+R/\sqrt{\Var I}$ with $r_n=\sqrt{2(n-1)/(2n+5)}\in[0,1]$ and $\E(R^2)/\Var I=7/(2n+5)=1-r_n^2$. Since $L$ and $R$ are
uncorrelated,
\[
\E(I'-L)^2=(1-r_n)^2+(1-r_n^2)=2(1-r_n)\le2(1-r_n^2)=\frac{14}{2n+5}.\qedhere
\]
\end{proof}

\section{The interaction graph}\label{sec:graph}

We use the normal approximation theorem of Chatterjee~\cite[Theorem~2.5]{Ch08} in the form stated by Chatterjee and
Diaconis~\cite[Theorem~4.1]{CD}. Let $\mathcal X$ be a measurable space and $\phi:\mathcal X^n\to\mathbb R$ measurable.
A graphical rule assigns to each $x\in\mathcal X^n$ a simple graph $\mathcal G(x)$ on $[n]$; it is symmetric if
relabelling the coordinates of $x$ relabels the graph accordingly. For $x,x^*\in\mathcal X^n$ let $x^i$ be $x$ with
$x_i$ replaced by $x_i^*$, and $x^{ij}$ the vector with both $x_i$ and $x_j$ replaced. The rule is an interaction
rule for $\phi$ if, whenever $\{i,j\}$ is an edge of none of $\mathcal G(x)$, $\mathcal G(x^i)$, $\mathcal G(x^j)$,
$\mathcal G(x^{ij})$, we have $\phi(x)-\phi(x^j)=\phi(x^i)-\phi(x^{ij})$. A symmetric graphical rule
$\mathcal G'$ on $\mathcal X^{n+4}$ is an extension of $\mathcal G$ if $\mathcal G(x_1,\dots,x_n)$ is a subgraph of
$\mathcal G'(x_1,\dots,x_{n+4})$ for all $x_1,\dots,x_{n+4}$.

\begin{theorem}[Chatterjee]\label{thm:chatterjee}
Let $\phi$ admit a symmetric interaction rule $\mathcal G$ with an extension $\mathcal G'$ on $\mathcal X^{n+4}$. Let
$X_1,X_2,\dots$ be independent and identically distributed in $\mathcal X$, $X=(X_1,\dots,X_n)$, $X^*$ an independent
copy of $X$, $\Phi=\phi(X)$ and $\sigma^2=\Var\Phi>0$. Let $\Delta_j\phi=\Phi-\phi(X^j)$ with $X^j$ formed from $X$
and $X^*$ as $x^j$ is formed from $x$ and $x^*$, $M=\max_j|\Delta_j\phi|$, and $\delta=1+$ the degree of the vertex $1$ in
$\mathcal G'(X_1,\dots,X_{n+4})$. Then, for a universal constant $C_0$,
\[
\dW\Bigl(\frac{\Phi-\E\Phi}{\sigma},Z\Bigr)\le\frac{C_0n^{1/2}}{\sigma^2}\bigl(\E M^8\bigr)^{1/4}\bigl(\E\delta^4\bigr)^{1/4}
+\frac1{2\sigma^3}\sum_{j=1}^n\E|\Delta_j\phi|^3.
\]
\end{theorem}

Take $\mathcal X=[0,1]^2$. Call $x\in\mathcal X^m$ generic if its first coordinates are distinct and its second
coordinates are distinct. For generic $x$ let $\pi_x$ be the permutation it induces as in Section~\ref{sec:proj}, and
for $a\in[m]$ let $\alpha_x(a)$ and $\beta_x(a)$ be the ranks of the first and second coordinate of $x_a$. For
non-generic $x$ define $\pi_x$ by breaking ties by the index, which makes every statistic below measurable; its values
there play no role.

\begin{lemma}\label{lem:graph}
Fix $k\ge2$, let $W(x)=F(\pi_x)+G(\pi_x)$ be a two-sided local statistic of window length $k$, let $s>0$, and let
$\phi=\lambda\sum_a\psi(x_a)+\mu W/s$ with real $\lambda,\mu$. Let $\mathcal G(x)$, for $x\in\mathcal X^n$, be the
complete graph if $x$ is not generic, and otherwise the graph in which $a\ne b$ are joined when
$|\alpha_x(a)-\alpha_x(b)|\le k-1$ or $|\beta_x(a)-\beta_x(b)|\le k-1$. Let $\mathcal G'$ on $\mathcal X^{n+4}$ be
defined in the same way with $k+3$ in place of $k-1$. Then $\mathcal G$ is a symmetric interaction rule for $\phi$,
$\mathcal G'$ is an extension of $\mathcal G$, every vertex of $\mathcal G'(x)$ has degree at most $4k+12$ when $x$ is
generic, and for generic $x$ and $x^j$
\[
|W(x)-W(x^j)|\le 8k-4.
\]
\end{lemma}

\begin{proof}
Symmetry is clear, since genericity and ranks are unchanged by relabelling. If $x_1,\dots,x_n$ is not generic, neither
is $x_1,\dots,x_{n+4}$, and both graphs are complete. If $x_1,\dots,x_{n+4}$ is not generic, $\mathcal G'$ is complete.
Otherwise, adding four points changes a difference of ranks among the first $n$ points by at most $4$, so every edge of
$\mathcal G(x_1,\dots,x_n)$ is an edge of $\mathcal G'(x_1,\dots,x_{n+4})$. In a generic configuration each vertex has
at most $2(k+3)$ neighbours in each order, which gives the degree bound.

For the interaction property, the part $\lambda\sum_a\psi(x_a)$ satisfies $\phi(x)-\phi(x^j)=\phi(x^i)-\phi(x^{ij})$
identically, because it is a sum of functions of one coordinate each. It therefore suffices to prove the identity for
$F(\pi_x)$; the argument for $G$ is the same with the roles of the two coordinates exchanged. Suppose $\{i,j\}$ is an
edge of none of the four graphs. Then the four configurations are generic, and in each of them the first-coordinate
ranks of the $i$-th and the $j$-th points differ by at least $k$.

Regard a generic configuration as the sequence of its points listed by first coordinate; $F$ adds $f$ over the
patterns, in the second coordinate, of all blocks of $k$ consecutive points. Removing a point $z$ deletes the blocks
that contain $z$ and creates the blocks that contain both the last points before $z$ and the first points after $z$;
every other block keeps its points and their relative order. So the change of $F$ caused by removing $z$ is a function
of the ordered list of $z$, the at most $k-1$ points immediately before $z$ and the at most $k-1$ points immediately
after it, through their second coordinates; it does not depend on the absolute position of $z$ in the sequence, which
is where the fact that $f$ is the same on all windows is used. Removing $z$ deletes at most $k$ blocks and creates
at most $k-1$, so this change is at most $2k-1$ in absolute value. The same holds for inserting a point.

Now $F(\pi_x)-F(\pi_{x^i})$ is the change caused by removing $x_i$ from $x$, plus the change caused by inserting
$x_i^*$ into the result. In $x$ there are at least $k-1$ points strictly between $x_i$ and $x_j$ in the first order, so
the $k-1$ points immediately before and after $x_i$ do not include $x_j$ and are not altered when $x_j$ is removed.
The same holds in $x^j$ for $x_i$ and $x_j^*$. Hence removing $x_i$ causes the same change in $x$ as in $x^j$, and the
neighbourhoods of $x_i$ in the two configurations, after $x_i$ is removed, still agree. Applying the same argument to
$x^i$ and $x^{ij}$ shows that inserting $x_i^*$ causes the same change in both. Therefore
$F(\pi_x)-F(\pi_{x^i})=F(\pi_{x^j})-F(\pi_{x^{ij}})$, which is the required identity. Finally, for generic $x$ and
$x^j$, $W(x)-W(x^j)$ is a removal and an insertion for $F$ and for $G$, at most $2(2k-1)$ each.
\end{proof}

\begin{remark}
The bound $|W(x)-W(x^j)|\le8k-4$ uses that $f$ and $g$ do not depend on the window's position. If different windows
carry different functions, removing a point shifts the positions of all later windows, and a single replacement can
change the statistic by an amount of order $n$.
\end{remark}

\section{Proofs}\label{sec:proofs}

\begin{proof}[Proof of Theorem~\ref{thm:main}]
Realize $\pi=\pi_X$ from independent uniform points $X=(P_1,\dots,P_n)$, let $s=\sqrt{\Var W}\ge\sqrt{cn}$, and fix
$a^2+b^2=1$. Apply Theorem~\ref{thm:chatterjee} to $\phi=6an^{-1/2}\sum_{l=1}^n\psi(x_l)+bW/s$, so that $\phi(X)-\E\phi(X)=aL+bW'$.
By Lemma~\ref{lem:graph} the hypotheses hold, and since $X$, $X^*$ and $X_1,\dots,X_{n+4}$ are generic with
probability one, $\delta\le4k+13$ and
\[
M\le\frac{6|a|}{\sqrt n}+\frac{(8k-4)|b|}{s}\le\frac{m_0}{\sqrt n},\qquad m_0=6+\frac{8k}{\sqrt c},
\]
with probability one, using $|\psi(p)-\psi(q)|\le1$. Next, $\sigma^2=\Var(aL+bW')=1+2ab\Cov(L,W')$. By Lemma~\ref{lem:cov},
$|\Cov(I',W')|\le2n\sqrt{k^3/24}/(\sqrt{\Var I}\,s)\le c_1n^{-1}$ for a constant $c_1$ depending on $k$ and $c$, and
by Lemma~\ref{lem:hoeff} and the Cauchy--Schwarz inequality
\[
|\Cov(L,W')|\le|\Cov(I',W')|+\|L-I'\|_2\le c_1n^{-1}+\sqrt{14/(2n+5)}.
\]
Here $\Var I=n(n-1)(2n+5)/72\ge n^3/36$ for $n\ge2$, so one may take $c_1=12\sqrt{k^3/24}/\sqrt c$. Choose
$n_1\ge\max(k,2)$ so that the right side is at most $1/4$ for $n\ge n_1$; then $\tfrac12\le\sigma^2\le\tfrac32$. Since
$(\E M^8)^{1/4}\le m_0^2/n$ and $(\E\delta^4)^{1/4}\le4k+13$, Theorem~\ref{thm:chatterjee} gives
\[
\varepsilon_n:=\dW\Bigl(\frac{aL+bW'}{\sigma},Z\Bigr)\le\frac{2C_0(4k+13)m_0^2}{\sqrt n}+\frac{n}{2\sigma^3}\Bigl(\frac{m_0}{\sqrt n}\Bigr)^3
\le\frac{2C_0(4k+13)m_0^2+\sqrt2\,m_0^3}{\sqrt n}.
\]
It remains to pass from $aL+bW'$ to $aI'+bW'$ and from $\sigma$ to $1$. The Wasserstein distance between two random
variables defined on one space is at most the mean of their absolute difference, and
$\dW(\sigma X,\sigma Y)=\sigma\,\dW(X,Y)$. Hence
\[
\dW(aI'+bW',Z)\le|a|\,\|I'-L\|_2+\sigma\varepsilon_n+|\sigma-1|\,\E|Z|.
\]
The first term is at most $\sqrt{14/(2n+5)}\le3n^{-1/2}$ by Lemma~\ref{lem:hoeff}. In the last,
$|\sigma-1|\le|\sigma^2-1|\le|\Cov(L,W')|\le c_1n^{-1}+3n^{-1/2}$. All three terms are at most a constant depending
on $k$ and $c$ times $n^{-1/2}$, which proves the bound.

For the convergence in distribution, $\liminf\Var(W_n)/n>0$ gives $c>0$ with $\Var W_n\ge cn$ for all large $n$, so
$aI'+bW_n'$ converges to $N(0,1)$, the law of $aZ_1+bZ_2$ for independent standard normal $Z_1,Z_2$, for every unit
vector $(a,b)$, and by scaling for every $(a,b)$. The Cramér--Wold device completes the proof.
\end{proof}

\begin{proof}[Proof of Corollary~\ref{cor:T}]
The covariance is Lemma~\ref{lem:cov} and the variance is computed in~\cite{CD}. Since $36\Var I/n^3\to1$ and
$6\Var T/n\to1$, the correlation $\frac{n-1}2\bigl(\Var I\Var T\bigr)^{-1/2}$ equals $(1+o(1))\frac n2\cdot\frac{6\sqrt6}{n^2}
=(3\sqrt6+o(1))/n$. The statistic $T$ is two-sided local with $k=2$, and $\Var T\ge n/6$ for all $n\ge2$, so
Theorem~\ref{thm:main} applies with $k=2$ and $c=1/6$; the constants it gives are universal.
\end{proof}

\begin{proof}[Proof of Corollary~\ref{cor:triple}]
Write $E=D(\pi^{-1})$ and $s_D^2=(n+1)/12$, the common variance of $D(\pi)$ and $E$. By Cramér--Wold it suffices to show
that $aI'+bD(\pi)'+cE'$ converges to $N(0,a^2+b^2+c^2)$ for every $(a,b,c)$. If $b=c=0$ this is the central limit
theorem for $I$, which follows from Lemma~\ref{lem:hoeff} and the central limit theorem for the independent sum $L$.
Otherwise let $m=\max(|b|,|c|)$ and $W=(b/m)D(\pi)+(c/m)E$, a two-sided local statistic of window length $2$ with
$f=(b/m)\mathbf 1\{\cdot=21\}$ and $g=(c/m)\mathbf 1\{\cdot=21\}$. From the variances recalled after
Lemma~\ref{lem:cov},
\[
m^2\Var W=(b^2+c^2)\frac{n+1}{12}+bc\,\frac{n-1}n\ge(b^2+c^2)\Bigl(\frac n{12}-\frac12\Bigr)\ge\frac{(b^2+c^2)n}{24}
\]
for $n\ge12$, and $b^2+c^2\ge m^2$, so $\Var W\ge n/24$. By Theorem~\ref{thm:main}, $(I',W')$ converges to a standard
normal vector in $\mathbb R^2$. Now $bD(\pi)'+cE'=\tau_nW'$ with $\tau_n=m\sqrt{\Var W}/s_D$, and
$\tau_n^2=b^2+c^2+12bc(n-1)/(n(n+1))\to b^2+c^2$. Hence $aI'+bD(\pi)'+cE'=aI'+\tau_nW'$ converges to
$N(0,a^2+b^2+c^2)$.
\end{proof}

\section*{Notation}

\noindent\begin{tabular}{p{0.3\textwidth}p{0.45\textwidth}p{0.15\textwidth}}
$[n]$, $S_n$, $\pi$ & the first $n$ positive integers; the symmetric group; a uniform element & Section~\ref{sec:intro}\\
$I$, $D$, $T$ & inversions; descents; $D(\pi)+D(\pi^{-1})$ & Section~\ref{sec:intro}\\
$k$, $\operatorname{st}$, $f$, $g$ & window length; pattern of a sequence; window functions & \eqref{eq:W}\\
$F$, $G$, $W$ & sums over the windows of $\pi$ and of $\pi^{-1}$; their sum & \eqref{eq:W}\\
$X'$, $Z$, $\dW$ & standardization; standard normal variable; Wasserstein distance & Section~\ref{sec:intro}\\
$\inv(\sigma)$ & inversions of $\sigma\in S_k$ & Section~\ref{sec:moments}\\
$P_a=(U_a,V_a)$, $h$ & uniform points of $[0,1]^2$; discordance kernel & Section~\ref{sec:proj}\\
$\psi$, $\tilde h$, $L$, $R$, $r_n$ & Hoeffding projection and remainder of $I$ & Lemma~\ref{lem:hoeff}\\
$\mathcal G$, $\mathcal G'$, $x^i$, $x^{ij}$ & interaction rule, its extension, replaced configurations &
Section~\ref{sec:graph}\\
$\Delta_j\phi$, $M$, $\delta$ & the quantities of Chatterjee's theorem & Section~\ref{sec:graph}\\
$\alpha_x$, $\beta_x$, $\pi_x$ & ranks of the coordinates; induced permutation & Section~\ref{sec:graph}\\
$c$, $C$, $C_0$, $c_1$, $m_0$, $n_1$ & constants & Section~\ref{sec:proofs}\\
\end{tabular}

\section*{Acknowledgements and funding}
This work was done by multiple AI agents, Claude (Anthropic) and Codex (OpenAI), connected to
Math Brain~\cite{MathBrainRepo} and working together inside the Brain framework~\cite{BrainFramework}, a shared,
persistent memory in which the agents record claims,
proofs, failed attempts and methods, retrieve one another's results, coordinate who works on what, and audit each
other's work. For more on the Brain and Math Brain, see the Math Brain paper~\cite{MathBrainPaper}, the Math Brain
repository~\cite{MathBrainRepo} and the Brain framework repository~\cite{BrainFramework}.

The research presented in this paper was supported by the European Research Council (ERC) under the European Union's
Horizon Europe research and innovation programme (grant agreement No.~101041711), by the Simons Foundation as part of
the Collaboration on the Mathematical and Scientific Foundations of Deep Learning, by Heights Labs, by Convex Nexus
Capital, by the Israel Science Foundation (grant number 2258/19), and by the Israel Science Foundation (ISF Grant
4101/25). Computational resources were provided in part by Google Cloud through its education credits programme
(credit 530550103).

\section*{Declaration of competing interest}
The author holds equity in Heights Labs and in Convex Nexus Capital and has a management role at Convex Nexus Capital.
These entities are included in the funding declaration above.

\section*{Data availability}
The code that checks Lemma~\ref{lem:cov} and the variance formulas by enumerating $S_n$ for $n\le8$ is in the
repository~\cite{Repo}.

\section*{Declaration of generative AI and AI-assisted technologies}
Claude agents (Anthropic) proved Theorem~\ref{thm:main} and its corollaries and prepared the manuscript; a Codex agent
(OpenAI) refereed the proofs. The coupling and the interaction graph are those of Chatterjee and Diaconis~\cite{CD}. The
question is Dörr's~\cite{Dorr}. The joint statements, which answer Dörr's question, are the new results; Lemma~\ref{lem:window} is the conditioning
identity they use. The author
directed the work and is responsible for the paper.

\begin{thebibliography}{99}
\bibitem{BruckRottger} B.~Brück and F.~Röttger, \emph{A central limit theorem for the two-sided descent statistic on
Coxeter groups}, Electron. J. Combin. \textbf{29} (2022), Paper P1.1, \arxiv{1908.07955}.
\bibitem{Ch08} S.~Chatterjee, \emph{A new method of normal approximation}, Ann. Probab. \textbf{36} (2008),
1584--1610, \doi{10.1214/07-AOP370}.
\bibitem{CD} S.~Chatterjee and P.~Diaconis, \emph{A central limit theorem for a new statistic on permutations}, Indian
J. Pure Appl. Math. \textbf{48} (2017), 561--573, \arxiv{1608.01666}.
\bibitem{Dorr} P.~Dörr, \emph{Joint asymptotic normality and extremes of generalized inversions and descents},
preprint, 2024; version~5, 2026, \arxiv{2404.06598}.
\bibitem{DorrHeiny} P.~Dörr and J.~Heiny, \emph{Joint extremes of inversions and descents of random permutations},
J. Theoret. Probab. \textbf{38} (2025), Paper No.~42, \arxiv{2309.17314}.
\bibitem{FangRollin} X.~Fang and A.~Röllin, \emph{Rates of convergence for multivariate normal approximation with
applications to dense graphs and doubly indexed permutation statistics}, Bernoulli \textbf{21} (2015), 2157--2189,
\doi{10.3150/14-BEJ639}.
\bibitem{FerayCoxeter} V.~Féray, \emph{On the central limit theorem for the two-sided descent statistics in Coxeter
groups}, Electron. Commun. Probab. \textbf{25} (2020), Paper No.~28, \arxiv{1911.10939}.
\bibitem{He} J.~He, \emph{A central limit theorem for descents of a Mallows permutation and its inverse}, Ann. Inst.
Henri Poincaré Probab. Stat. \textbf{58} (2022), 667--694, \arxiv{2005.09802}.
\bibitem{Hoeffding} W.~Hoeffding, \emph{A class of statistics with asymptotically normal distribution}, Ann. Math.
Statist. \textbf{19} (1948), 293--325, \doi{10.1214/aoms/1177730196}.
\bibitem{Kendall} M.~G.~Kendall, \emph{A new measure of rank correlation}, Biometrika \textbf{30} (1938), 81--93,
\doi{10.1093/biomet/30.1-2.81}.
\bibitem{Rottger} F.~Röttger, \emph{Asymptotics of a locally dependent statistic on finite reflection groups},
Electron. J. Combin. \textbf{27} (2020), Paper P2.24, \arxiv{1812.00372}.
\bibitem{Vatutin} V.~A.~Vatutin, \emph{The numbers of ascending segments in a random permutation and in one inverse to
it are asymptotically independent}, Discrete Math. Appl. \textbf{6} (1996), 41--52, \doi{10.1515/dma.1996.6.1.41}.
\bibitem{Repo} Y.~Shmalo, \emph{Inversions and two-sided descents: paper and enumeration program}, 2026,
\url{https://github.com/yspennstate/inversions-two-sided-descents-clt}.
\bibitem{MathBrainRepo} Y.~Shmalo, \emph{Math Brain: public snapshot, code and tests}, 2026,
\url{https://github.com/yspennstate/math-brain-public/tree/8c9a6bd478c4ec4082347151b7ae16db50c917dd}.
\bibitem{MathBrainPaper} Y.~Shmalo, \emph{Math Brain: shared, auditable memory for multi-agent mathematical research},
2026, \doi{10.5281/zenodo.24150837},
\url{https://github.com/yspennstate/math-brain-public/tree/main/papers/math-brain-2026-10-24}.
\bibitem{BrainFramework} Y.~Shmalo, \emph{Brain Framework}, public software repository, 2026,
\url{https://github.com/yspennstate/brain-framework-public/tree/5629e644cdd510173ae97e9aa92642dabd22dc57}.
\end{thebibliography}

\end{document}
